% A Model of Compounding — standalone preview (compiles with pdflatex).
% AWR loss + BoN variance reduction + the compounding model.
% NOTE: present as a MODEL (first-order approximations), not theorems.
\documentclass[11pt]{article}
\usepackage[margin=1in]{geometry}
\usepackage{amsmath,amssymb}
\usepackage{bm}
\newcommand{\E}{\mathbb{E}}
\newcommand{\KL}{\mathrm{KL}}
\newcommand{\R}{R}
\title{\vspace{-2em}A Model of Compounding\\ \large (preview: AWR loss, BoN variance reduction)}
\date{}
\begin{document}
\maketitle

\paragraph{Setup and notation.}
A closed-loop action-chunking policy $\pi_\theta$ acts by, at each replan $k$, observing state
$s_k$ and sampling an \emph{ordered action chunk} $a=z_{1:K}\sim\pi_\theta(\cdot\mid s_k)$ of $K$
tokens, executing its first $R$ steps, then replanning. An episode of horizon $T$ therefore
contains
\begin{equation}
  H \;=\; \big\lceil T/R \big\rceil
  \label{eq:H}
\end{equation}
replans. Write the (sparse) episode return as $\R(s,a)\in\{0,1\}$ (task success), and let
$V(s)=\E_{\pi_\theta}[\R\mid s]$ be the state value and $A(s,a)=\R(s,a)-V(s)$ the advantage. The
per-token factorization is $\log\pi_\theta(a\mid s)=\sum_{k=1}^{K}\log\pi_\theta(z_k\mid z_{<k},s)$.

%======================================================================
\section*{1. Per-observation adaptivity has no exploitable signal}
For a per-observation compute choice $c$ (token depth $K$ or replan horizon $R$), define its
\emph{single-deviation oracle value} by counterfactual rollout (execute $c$ at $s$, then follow
the default policy):
\begin{equation}
  \Delta(s,c)\;=\;Q(s,c)-Q(s,c_{\text{def}}),\qquad
  \E_{s}\big[\,\Delta(s,c)\,\big]\;\approx\;0 .
  \label{eq:oracle}
\end{equation}
Equivalently, the per-step failure rate $\rho(s,c)$ (probability the chunk is
\emph{unrecoverable}) is (nearly) independent of $c$: closed-loop replanning corrects a single
cheap decision, so $\rho(s,c)\approx\rho(s)$. There is thus no per-state ``where to economize''
signal to exploit.

%======================================================================
\section*{2. Selection as mode-seeking / tail-risk removal (best-of-$N$)}
At a fixed budget, draw $N$ candidate chunks $a_1,\dots,a_N\stackrel{\text{iid}}{\sim}\pi_\theta(\cdot\mid s)$
and execute the \emph{consensus} candidate (kernel-density vote):
\begin{equation}
  \hat a \;=\; \arg\max_{i\in[N]}\ \sum_{j=1}^{N}\kappa\!\left(a_i,a_j\right)
  \;\xrightarrow{\,N\to\infty\,}\; \arg\max_{a}\ \pi_\theta(a\mid s),
  \label{eq:vote}
\end{equation}
i.e.\ consensus is a \emph{mode-seeking} operator: it converts sampling from $\pi_\theta$ into
(approximate) selection of its mode. It does \emph{not} seek the ``best'' chunk (by
\eqref{eq:oracle} there is none) and it does \emph{not} average; it \emph{rejects low-density
outliers}.

Its benefit is asymmetric, not symmetric noise reduction. Model a sampled chunk as
\emph{catastrophic} (unrecoverable, outcome $u_-$) with probability $\rho$, else safe (outcome
$u_+>u_-$); the mode is safe. The single-sample outcome averages the bad tail,
$\E_{a\sim\pi_\theta}[u]=(1-\rho)u_+ + \rho\,u_-$, whereas the consensus removes it, giving a
per-step gain equal to the removed tail risk,
\begin{equation}
  g \;=\; u(\hat a)-\E_{a\sim\pi_\theta}[u]\;=\;\rho\,(u_+-u_-)\;=\;\rho\,\Delta .
  \label{eq:gain}
\end{equation}
Equivalently, the per-step \emph{failure} rate drops from $\rho$ to
\begin{equation}
  \rho_N \;=\; \Pr\big[\hat a \text{ catastrophic}\big]
        \;\approx\; \Pr\!\big[\mathrm{Binom}(N,\rho)\ge\lceil N/2\rceil\big]
        \;\approx\; \binom{N}{\lceil N/2\rceil}\,\rho^{\lceil N/2\rceil}
        \;\ll\; \rho .
  \label{eq:rhoN}
\end{equation}
Thus best-of-$N$ improves outcomes by pruning the rare catastrophic tail
($\rho\!\to\!\rho_N$), not by averaging and not by ranking quality; the gain \eqref{eq:gain}
exists precisely because the tail is asymmetrically bad.

%======================================================================
\section*{3. The compounding law}
Assuming episode success requires avoiding an unrecoverable chunk at every replan, the
catastrophe-free success factor is $\prod_{k=1}^{H}(1-\rho)\approx(1-\rho)^{H}$. Comparing a
single sample ($\rho$) with best-of-$N$ ($\rho_N$):
\begin{equation}
  \Delta\mathrm{SR}
  \;=\;(1-\rho_N)^{H}-(1-\rho)^{H}
  \;\approx\; H\,(\rho-\rho_N)\qquad(\rho,\rho_N\ \text{small}).
  \label{eq:compound}
\end{equation}
Equation~\eqref{eq:compound} is the \textbf{compounding law}: the tiny per-step benefit
$(\rho-\rho_N)$ is multiplied by the number of replans $H$, because a \emph{single} catastrophe
anywhere fails the episode. It yields a falsifiable prediction,
\begin{equation}
  \boxed{\ \Delta\mathrm{SR}\ \propto\ H\;=\;\#\text{replans}\ \times\ (\rho-\rho_N)\ }
  \label{eq:predict}
\end{equation}
i.e.\ the selection gain grows with episode length / replan frequency and with per-step
head-room ($\rho$). \emph{The same $H$ factor explains the asymmetry with adaptivity:} by
\eqref{eq:oracle} a budget choice leaves $\rho$ unchanged, so it multiplies zero
($\Delta\mathrm{SR}\approx0$), whereas selection reduces $\rho$, so it compounds.

%======================================================================
\section*{4. Distilling selection: advantage-weighted regression (AWR)}
Best-of-$N$ shifts the executed-action distribution toward the consensus mode. We bake this
shift into a single-forward policy by advantage-weighted regression. AWR solves the
KL-regularized improvement
\begin{equation}
  \max_{\pi}\ \E_{a\sim\pi}\!\big[A(s,a)\big]-\beta\,\KL\!\big(\pi(\cdot\mid s)\,\|\,\pi_{\!D}(\cdot\mid s)\big)
  \ \Longrightarrow\
  \pi^\star(a\mid s)\ \propto\ \pi_{\!D}(a\mid s)\,\exp\!\big(A(s,a)/\beta\big),
  \label{eq:awr-sol}
\end{equation}
and projects $\pi^\star$ onto $\pi_\theta$ by weighted maximum likelihood. With a frozen
reference $\pi_{\text{ref}}$ (anti-collapse) and the per-token factorization, the training loss is
\begin{equation}
  \mathcal{L}_{\text{AWR}}(\theta)
  = -\,\E_{(s,a)\sim\mathcal{D}}
     \!\left[\, w(s,a)\sum_{k=1}^{K}\log\pi_\theta(z_k\mid z_{<k},s)\right]
   \;+\;\beta_{\KL}\,\E_{s}\!\big[\KL\!\big(\pi_\theta(\cdot\mid s)\,\|\,\pi_{\text{ref}}(\cdot\mid s)\big)\big],
  \label{eq:awr-loss}
\end{equation}
\begin{equation}
  w(s,a)=\min\!\Big(\exp\!\big(A(s,a)/\beta\big),\,w_{\max}\Big),\qquad
  A(s,a)=\R(s,a)-V(s),
  \label{eq:awr-weight}
\end{equation}
with weights normalized to unit mean. Because $w\ge 0$, AWR can only \emph{raise} the likelihood
of observed chunks: it is imitation-bounded (it reproduces the selected/consensus behavior at a
single forward pass, but cannot exceed the source), consistent with our empirical ceiling.

%======================================================================
\section*{5. Cost model (why selection is nearly free)}
Perception dominates the per-replan cost, $C_{\text{vis}}\gg c_{\text{AR}}$, and is computed once
per replan and shared across the $N$ candidates:
\begin{equation}
  C_{\text{ep}}^{\text{single}} = H\,(C_{\text{vis}}+K\,c_{\text{AR}}),\qquad
  C_{\text{ep}}^{\text{BoN-}N} = H\,(C_{\text{vis}}+N K\,c_{\text{AR}})
  \ \xrightarrow[C_{\text{vis}}\gg c_{\text{AR}}]{}\ H\,C_{\text{vis}} .
  \label{eq:cost}
\end{equation}
Thus best-of-$N$ adds cost only on the cheap autoregressive axis; distillation
\eqref{eq:awr-loss} removes even that, retaining the gain at single-forward cost.

\bigskip
\noindent\textit{Caveat.} The safe/catastrophic dichotomy and the small-$\rho$ expansions are a
first-order \emph{model}, not a theorem; it is intended to explain the direction of the effects
and the prediction~\eqref{eq:predict}, which we test empirically (selection gain vs.\ \#replans).

\end{document}
